\documentclass[11pt,letterpaper]{article}
\PassOptionsToPackage{hyphens}{url}
\usepackage[margin=.83in,top=.8in,bottom=.8in,headheight=15pt]{geometry}
\usepackage[T1]{fontenc}
\usepackage{lmodern,amsmath,amssymb,mathtools,bm,booktabs,array,graphicx,tikz}
\usetikzlibrary{arrows.meta,positioning,calc,decorations.pathreplacing,patterns}
\definecolor{ink}{HTML}{17344D}\definecolor{accent}{HTML}{087E8B}\definecolor{pale}{HTML}{EDF5F5}\definecolor{gold}{HTML}{AD6B19}
\providecommand{\dd}{\,\mathrm{d}}\providecommand{\R}{\mathbb{R}}
\providecommand{\E}{\mathbb{E}}\providecommand{\Var}{\operatorname{Var}}
\providecommand{\Cov}{\operatorname{Cov}}\providecommand{\diag}{\operatorname{diag}}
\providecommand{\norm}[1]{\left\lVert#1\right\rVert}
\providecommand{\abs}[1]{\left|#1\right|}
\usepackage{microtype,hyperref,fancyhdr,enumitem,amsthm,sectsty,needspace}
\allsectionsfont{\raggedright}
\hypersetup{colorlinks=true,linkcolor=ink,urlcolor=accent,pdfauthor={Zan Ahmad},pdftitle={A reduced cardiac reaction--diffusion wave / Review}}
\setlength{\parindent}{0pt}\setlength{\parskip}{6pt}\linespread{1.07}
\setlist{itemsep=3pt,topsep=4pt,leftmargin=*}
\emergencystretch=2em\clubpenalty=10000\widowpenalty=10000
\pagestyle{fancy}\fancyhf{}\fancyhead[L]{\small Computational Cardiology}\fancyhead[R]{\small Lecture 10}
\fancyfoot[L]{\footnotesize Zan Ahmad}\fancyfoot[R]{\footnotesize\thepage}\renewcommand{\headrulewidth}{.3pt}
\newtheorem{proposition}{Proposition}\newtheorem{theorem}{Theorem}\newtheorem{lemma}{Lemma}
\theoremstyle{definition}\newtheorem{definition}{Definition}\newtheorem{example}{Example}
\numberwithin{equation}{section}
\begin{document}
\section*{A reduced cardiac reaction--diffusion wave --- review sheet}
Use this after attempting the worked examples and practice. Explain each claim in words before calculating.

\section*{Key ideas}
\begin{itemize}
\item The reduced cardiac model uses scaled voltage and a recovery gate, not the neuronal HH state.
\item Shared-face diffusion cancels internally under sealed boundaries.
\item Only diffusion conserves the stated voltage integral; reaction and stimulus change it.
\item Physical stimulus overlap and interpolated sensor coordinates keep refinement experiments consistent.
\item Component verification and active-wave sensitivity answer different questions.
\end{itemize}
\section*{Common errors and repairs}
\textbf{Reporting scaled voltage as millivolts without a conversion.} Keep the dimensionless axis or explicitly state and justify a voltage scale.

\textbf{Treating a missing arrival as zero speed.} Report the observation window and undefined two-sensor speed.

\textbf{Changing both mesh and time step and assigning an error order from one table.} Separate the controls and use an appropriate reference or qualified successive-difference comparison.

\section*{Retrieve without your notes}
\begin{itemize}
\item What are the units of D here?
\item Why can backward-Euler diffusion preserve constants?
\item What is one cm/ms in m/s?
\item What extra evidence is needed for biological validation?
\end{itemize}
\section*{Readiness check}
\begin{itemize}
\item Write the scaled two-current cardiac reaction equations with explicit units and limitations.
\item Derive sealed-strand diffusion as a shared-face flux balance.
\item Implement a threshold-crossing interpolation and convert a two-sensor delay into speed.
\item Separate rest, diffusion conservation, gate admissibility and numerical accuracy checks.
\item Compare spatial and temporal refinement at fixed physical sensor and stimulus locations.
\end{itemize}
\end{document}
